\documentclass[12pt]{article}
\usepackage[a4paper,margin=1in]{geometry}
\usepackage[utf8]{inputenc}
\usepackage[T1]{fontenc}
\usepackage{amsmath,amssymb}
\usepackage{textcomp}
\usepackage{parskip}
\begin{document}
\begin{center}
{\Large \textbf{2012 STEP I}}
\end{center}

\subsection*{Page 1}
\begin{flushleft}
\phi\par
\bigskip\textbf{Section A: Pure Mathematics}\par
\bigskip\textbf{Question 1.} The line L has equation y = c mx, with m > 0 and c > 0. It passes through the point\par
-\par
R(a, b) and cuts the axes at the points P (p, 0) and Q(0, q), where a, b, p and q are all positive.\par
Find p and q in terms of a, b and m.\par
As L varies with R remaining fixed, show that the minimum value of the sum of the distances\par
of P and Q from the origin is (a 2 1 + b 1 2 )2, and find in a similar form the minimum distance\par
between P and Q. (You may assume that any stationary values of these distances are minima.)\par
\bigskip\textbf{Question 2.} (i) Sketch the curve y = x4 6x2 + 9 giving the coordinates of the stationary points.\par
-\par
Let n be the number of distinct real values of x for which\par
4 2\par
\[x 6x + b = 0.\]
-\par
\[State the values of b, if any, for which (a) n = 0 ; (b) n = 1 ; (c) n = 2 ; (d) n = 3 ;\]
\[(e) n = 4 .\]
\textbf{(ii)} For which values of a does the curve y = x4 6x2 + ax + b have a point at which both\par
\[dy d2y -\]
\[= 0 and = 0 ?\]
\[dx dx2\]
For these values of a, find the number of distinct real values of x for which\par
4 2\par
\[x 6x + ax + b = 0 ,\]
-\par
in the different cases that arise according to the value of b.\par
\textbf{(iii)} Sketch the curve y = x4 6x2 + ax in the case a > 8 .\par
-\par
\end{flushleft}
\subsection*{Page 2}
\begin{flushleft}
\[\pi\]
\bigskip\textbf{Question 3.} (i) Sketch the curve y = sin x for 0 \leq x \leq 1 \pi and add to your diagram the tangent to\par
the curve at the origin and the chord joining the origin to the point (b, sin b), where\par
\[0 < b < \pi.\]
By considering areas, show that\par
1 2 1\par
\[1 b < cos b < 1 b sin b .\]
- 2 - 2\par
x\par
\textbf{(ii)} By considering the curve y = a , where a > 1, show that\par
2(a 1)\par
\[- < ln a < 1 + \sqrt\{2a\} 1 .\]
a + 1 - -\par
[Hint: You may wish to write a x as e x ln a .]\par
1 1\par
\bigskip\textbf{Question 4.} The curve C has equation xy = . The tangents to C at the distinct points P p, and\par
2 2p\par
Q q, , where p and q are positive, intersect at T and the normals to C at t\leqhese p\geqoints\par
2q\par
intersect at N . Show that T is the point\par
\[\leq \geq\]
2pq 1\par
, .\par
p + q p + q\par
( )\par
\[In the case pq = , find the coordinates of N . Show (in this case) that T and N lie on the\]
\[line y = x and are such that the product of their distances from the origin is constant.\]
\bigskip\textbf{Question 5.} Show that\par
\[1 \pi\]
4 1\par
\[sin(2x) ln(cos x) dx = (ln 2 1) ,\]
4 -\par
\[\int\]
and that\par
\[1 \pi\]
4 1\par
\[cos(2x) ln(cos x) dx = (\pi ln 4 2) .\]
8 - -\par
\[\int\]
Hence evaluate\par
\[1 \pi\]
cos(2x) + sin(2x) ln cos x + sin x dx .\par
\[1 \pi\]
\[\int 4\]
\[\leq \geq \leq \geq\]
\end{flushleft}
\subsection*{Page 3}
\begin{flushleft}
\sigma\par
\bigskip\textbf{Question 6.} A thin circular path with diameter AB is laid on horizontal ground. A vertical flagpole is\par
erected with its base at a point D on the diameter AB. The angles of elevation of the top\par
of the flagpole from A and B are α and β respectively (both are acute). The point C lies on\par
the circular path with DC perpendicular to AB and the angle of elevation of the top of the\par
\[flagpole from C is φ. Show that cot α cot β = cot2 φ.\]
Show that, for any p and q,\par
2 1 2 1 1 1\par
\[cos p cos q sin (p + q) sin p sin q cos (p + q) = cos(p + q) cos(p + q) cos(p q).\]
2 - 2 2 - 2 -\par
\[Deduce that, if p and q are positive and p + q \leq 1 \pi, then\]
\[cot p cot q \geq cot 2 1 (p + q)\]
\[and hence show that φ \leq 1 (α + β) when α + β \leq 1 \pi .\]
2 2\par
\bigskip\textbf{Question 7.} A sequence of numbers t0, t1, t2, . . . satisfies\par
\[tn+2 = ptn+1 + qtn (n \geq 0),\]
where p and q are real. Throughout this question, x, y and z are non-zero real num-\par
bers.\par
\textbf{(i)} Show that, if tn = x for all values of n, then p + q = 1 and x can be any (non-zero) real\par
number.\par
\textbf{(ii)} Show that, if t2n = x and t2n+1 = y for all values of n, then q p = 1. Deduce that\par
±\par
\[either x = y or x = y, unless p and q take certain values that you should identify.\]
-\par
\textbf{(iii)} Show that, if t3n = x, t3n+1 = y and t3n+2 = z for all values of n, then\par
3 3\par
\[p + q + 3pq 1 = 0 .\]
-\par
\[Deduce that either p + q = 1 or (p q)2 + (p + 1)2 + (q + 1)2 = 0. Hence show that\]
-\par
\[either x = y = z or x + y + z = 0.\]
\end{flushleft}
\subsection*{Page 4}
\begin{flushleft}
\omega\par
\bigskip\textbf{Question 8.} (i) Show that substituting y = xv, where v is a function of x, in the differential equation\par
dy\par
2 2\par
\[xy + y 2x = 0 (x = 0)\]
\[dx - \geq\]
leads to the differential equation\par
dv\par
\[xv + 2v 2 = 0 .\]
dx -\par
Hence show that the general solution can be written in the form\par
2 2 2\par
\[x (y x ) = C ,\]
-\par
where C is a constant.\par
\textbf{(ii)} Find the general solution of the differential equation\par
dy\par
\[y + 6x + 5y = 0 (x = 0).\]
\[dx \geq\]
\end{flushleft}
\subsection*{Page 5}
\begin{flushleft}
)\par
\bigskip\textbf{Section B: Mechanics}\par
\bigskip\textbf{Question 9.} A tall shot-putter projects a small shot from a point 2.5 m above the ground, which is horizon-\par
tal. The speed of projection is 10 m s\par
-1\par
and the angle of projection is θ above the horizontal.\par
Taking the acceleration due to gravity to be 10 m s\par
-2,\par
show that the time, in seconds, that\par
elapses before the shot hits the ground is\par
\[\sqrt\{1\} c + \sqrt\{2\} c ,\]
\[\sqrt\{2\} - -\]
\[\leq \geq\]
\[where c = cos 2θ.\]
\[Find an expression for the range in terms of c and show that it is greatest when c = .\]
Show that the extra distance attained by projecting the shot at this angle rather than at an\par
\[angle of 45 ^\circ is 5(\sqrt\{6\} \sqrt\{2\} 1) m.\]
- -\par
\bigskip\textbf{Question 10.} I stand at the top of a vertical well. The depth of the well, from the top to the surface of\par
the water, is D. I drop a stone from the top of the well and measure the time that elapses\par
between the release of the stone and the moment when I hear the splash of the stone entering\par
the water.\par
In order to gauge the depth of the well, I climb a distance δ down into the well and drop a\par
stone from my new position. The time until I hear the splash is t less than the previous time.\par
Show that\par
2D 2(D δ) δ\par
\[t = - + ,\]
\[\sum g - \sum g u\]
where u is the (constant) speed of sound. Hence show that\par
1 2\par
\[D = gT ,\]
1 δ δ\par
\[where T = β + and β = t .\]
2 βg - u\par
\[Taking u = 300 m s\]
-1\par
\[and g = 10 m s\]
-2,\par
\[show that if t =\]
\[s and δ = 10 m, the well is\]
approximately 185 m deep.\par
\end{flushleft}
\subsection*{Page 6}
\begin{flushleft}
(\par
\bigskip\textbf{Question 11.} The diagram shows two particles, A of mass 5m and B of mass 3m, connected by a light\par
inextensible string which passes over two smooth, light, fixed pulleys, Q and R, and under a\par
smooth pulley P which has mass M and is free to move vertically.\par
Particles A and B lie on fixed rough planes inclined to the horizontal at angles of arctan and\par
arctan respectively. The segments AQ and RB of the string are parallel to their respective\par
planes, and segments QP and P R are vertical. The coefficient of friction between each particle\par
and its plane is μ.\par
Q R\par
A\par
B\par
P\par
\textbf{(i)} Given that the system is in equilibrium, with both A and B on the point of moving up\par
\[their planes, determine the value of μ and show that M = 6m.\]
\textbf{(ii)} In the case when M = 9m, determine the initial accelerations of A, B and P in terms\par
of g.\par
\end{flushleft}
\subsection*{Page 7}
\begin{flushleft}
)\par
\bigskip\textbf{Section C: Probability and Statistics}\par
12 Fire extinguishers may become faulty at any time after manufacture and are tested annually\par
on the anniversary of manufacture.\par
The time T years after manufacture until a fire extinguisher becomes faulty is modelled by\par
the continuous probability density function\par
2t\par
\[for t \geq 0 ,\]
\[(1 + t^2)2\]
\[f(t) =\]
⎧\par
⎪ ⎨ 0 otherwise.\par
⎪\par
A faulty fire extinguisher will fail an an⎩nual test with probability p, in which case it is destroyed\par
immediately. A non-faulty fire extinguisher will always pass the test. All of the annual tests\par
are independent.\par
Show that the probability that a randomly chosen fire extinguisher will be destroyed exactly\par
\[three years after its manufacture is p(5p2 13p + 9)/10.\]
-\par
Find the probability that a randomly chosen fire extinguisher that was destroyed exactly three\par
years after its manufacture was faulty 18 months after its manufacture.\par
\bigskip\textbf{Question 13.} I choose at random an integer in the range 10000 to 99999, all choices being equally likely.\par
Given that my choice does not contain the digits 0, 6, 7, 8 or 9, show that the expected number\par
of different digits in my choice is 3.3616.\par
\end{flushleft}
\end{document}
